\documentclass{article}

\usepackage{float}
\restylefloat{table}

\usepackage{booktabs}

\title{Team Productivity: POC\\\progname}

\author{\authname}

\date{}

\input{../Comments.text}
\input{../Common.text}

\begin{document}

\maketitle

This document summarizes the contributions of each team member up to the POC
Demo.  The time period of interest is September 28, 2026 through November 20,
2026.

\wss{Please do not delete sections or modify the table formats.  Standardization
helps the instructors and TAs to review these documents. You may add sections to
the report.}

\wss{Remove the instructor guidelines/comments in the final version.}

\wss{Fill in the requested information, like project url, TA name, supervisor name, etc.}

\section{Demo Plans}

The primary feasibility risks are mobile audio-lyric synchronization,
language-specific lyric segmentation, and delivering flashcards to Anki from
a phone.  Mobile playback must maintain line highlighting within 0.5 seconds,
Japanese and mixed-script lyric text must be segmented reliably despite song
slang and fragments, and the team must determine whether AnkiConnect,
AnkiDroid, or an export file is the practical integration method.

The PoC will use one test phone and one short test track per demonstrated
language.  We will demonstrate:
\begin{itemize}
    \item real-time, line-by-line lyric highlighting synchronized to audio,
    with at least 95\% of line changes occurring within 0.5 seconds of the
    mapped timestamp;
    \item selectable, romanized Japanese words from 20 checked lyric lines,
    with at least 90\% correct segmentation and text ready within three
    seconds; and
    \item creation of a flashcard containing the selected term, its translation,
    and context sentence in a test Anki deck.
\end{itemize}

Each component owner will demonstrate their corresponding risk.  Omar Bakr will
cover mobile playback and learner-app behavior, Rhea Gokhale will cover
language segmentation, and Aansh Kotian will cover Anki/flashcard handoff.
Raymond Ma will act as the QA reviewer and will verify the test songs and
acceptance measurements.  The phone display will be mirrored to a laptop
screen, and the team will use a laptop for the presentation of logs and
diagnostics.  The demonstration will not include a slide deck or finished
user-interface styling.

If a live demonstration fails, the team will retain recorded logs and a
working test fixture for fallback evidence.  If synchronization is unreliable,
the team will use a coarser highlighting tolerance or a native audio module;
if segmentation is unreliable on the phone, segmentation will move into the
creator/editor stage for verification; and if direct Anki sending is blocked,
the team will fall back to an importable Anki export file.

\section{Team Meeting Attendance}

\wss{For each team member how many team meetings have they attended over the
time period of interest.  This number should be determined from the meeting
issues in the team's repo.  The first entry in the table should be the total
number of team meetings held by the team.}

\begin{table}[H]
\centering
\begin{tabular}{ll}
\toprule
\textbf{Student} & \textbf{Meetings}\\
\midrule
Total & Num\\
Name 1 & Num\\
Name 2 & Num\\
Name 3 & Num\\
Name 4 & Num\\
Name 5 & Num\\
\bottomrule
\end{tabular}
\end{table}

\wss{If needed, an explanation for the counts can be provided here.}

\section{Supervisor/Stakeholder Meeting Attendance}

\wss{For each team member how many supervisor/stakeholder team meetings have
they attended over the time period of interest.  This number should be determined
from the supervisor meeting issues in the team's repo.  The first entry in the
table should be the total number of supervisor and team meetings held by the
team.  If there is no supervisor, there will usually be meetings with
stakeholders (potential users) that can serve a similar purpose.}

\noindent \textbf{Supervisor's Name: } [fill in this information]

\begin{table}[H]
\centering
\begin{tabular}{ll}
\toprule
\textbf{Student} & \textbf{Meetings}\\
\midrule
Total & Num\\
Name 1 & Num\\
Name 2 & Num\\
Name 3 & Num\\
Name 4 & Num\\
Name 5 & Num\\
\bottomrule
\end{tabular}
\end{table}

\wss{If needed, an explanation for the counts can be provided here.}

\section{Lecture Attendance}

\wss{For each team member how many lectures have they attended over the time
period of interest.  This number should be determined from the lecture issues in
the team's repo. You can find the number of lectures in the time period of
interest by looking at the
\href{https://calendar.google.com/calendar/u/0/embed?src=rnboqiaki1k2la7rpu3bn0um58@group.calendar.google.com&ctz=America/Toronto}
{Google calendar} for the capstone course. The first entry in the
table should be the total number of lectures held for the class during the period of interest.}

\wss{NOTE: There will be approximately 13 lectures between the start of class
and the POC demos}

\begin{table}[H]
\centering
\begin{tabular}{ll}
\toprule
\textbf{Student} & \textbf{Lectures}\\
\midrule
Total & Num\\
Name 1 & Num\\
Name 2 & Num\\
Name 3 & Num\\
Name 4 & Num\\
Name 5 & Num\\
\bottomrule
\end{tabular}
\end{table}

\wss{If needed, an explanation for the lecture attendance can be provided here.}

\section{TA Document Discussion Attendance}

\wss{For each team member how many of the informal document discussion meetings
with the TA were attended over the time period of interest. The first row is the
total number of discussion meetings scheduled by the TA with the team.}

\noindent \textbf{TA's Name: } [fill in this information]

\begin{table}[H]
\centering
\begin{tabular}{ll}
\toprule
\textbf{Student} & \textbf{Meetings}\\
\midrule
Total & Num\\
Name 1 & Num\\
Name 2 & Num\\
Name 3 & Num\\
Name 4 & Num\\
Name 5 & Num\\
\bottomrule
\end{tabular}
\end{table}

\wss{If needed, an explanation for the attendance can be provided here.}

\section{Commits, Lines Added and Lines Deleted, Pull Requests}

\wss{For each team member how many commits to the main branch have been made
over the time period of interest.  The total is the total number of commits for
the entire team since the beginning of the term.  The percentage is the
percentage of the total commits made by each team member.}

\wss{There are also columns for lines added and lines deleted to give a sense of
the magnitude of each team member's contributions.  Note that these numbers can
be somewhat misleading, as a team member who does refactoring may have high
numbers in both columns, while a team member who implements a large new feature
may have a high number of lines added, but few lines deleted.  As for all commit
information, these numbers will be interpreted with care.}

\wss{The pull requests column counts the number of pull requests created by each
team member.  This is a measure of how many contributions were made in branches,
which may not be reflected in the commit counts for the main branch.}

\begin{table}[H]
\centering
\begin{tabular}{lllllll}
\toprule
\textbf{Student} & \textbf{Commits} & \textbf{Percent} & \textbf{Lines Added} &
\textbf{Lines Deleted} & \textbf{Pull Requests}\\
\midrule
Total & Num & 100\% & total added & total deleted & total PRs\\
Aansh Kotian & 12 & \% & 2,031 & 829 & 5\\
Name 2 & Num & \% & num & num & num\\
Name 3 & Num & \% & num & num & num\\
Name 4 & Num & \% & num & num & num\\
Name 5 & Num & \% & num & num & num\\
\bottomrule
\end{tabular}
\end{table}

\wss{If needed, an explanation for the counts can be provided here.  For
instance, if a team member has more commits to unmerged branches, these numbers
can be provided here.  If multiple people contribute to a commit, git allows for
multi-author commits.}

\section{Issue Tracker}

\wss{For each team member how many issues have they authored (including open and
closed issues (O+C)) and how many have they been assigned (only counting closed
issues (C only)) over the time period of interest.}

\begin{table}[H]
\centering
\begin{tabular}{lll}
\toprule
\textbf{Student} & \textbf{Authored (O+C)} & \textbf{Assigned (C only)}\\
\midrule
Aansh Kotian & 4 & 5 \\
Name 2 & Num & Num \\
Name 3 & Num & Num \\
Name 4 & Num & Num \\
Name 5 & Num & Num \\
\bottomrule
\end{tabular}
\end{table}

\wss{If needed, an explanation for the counts can be provided here.}

\section{CICD via GitHub Actions}

GitHub Actions is used to validate pull requests and keep the published project
documentation synchronized with the source.  The
\href{https://github.com/rhea333/LingoLyrics/blob/main/.github/workflows/latex-pages.yml}
{\texttt{latex-pages.yml}} workflow runs for pushes and pull requests that
change documents or references, builds changed LaTeX and Markdown files, and
deploys the compiled PDFs to GitHub Pages on the \texttt{main} branch.
Pull requests receive the build checks but do not publish the generated site.
Pushes to \texttt{main} also commit the generated PDFs back to the repository
and publish the site.

The Development Plan also specifies code CI for later implementation:
\texttt{ESLint}, \texttt{Prettier}, \texttt{Jest} coverage, and JSON
song-map schema validation.  Those checks will run on pull requests and pushes
to \texttt{main}, with pull requests required to pass CI and receive an
approved review before merging.

\section{Custom Documents (CDs)}

According to the team's
\href{https://gitlab.cas.mcmaster.ca/courses/capstone/-/blob/main/Teams/TeamComposition.csv}
{Team Composition entry}, LingoLyrics plans to produce a
\textbf{User/Stakeholder/Client Interview Report} in the fall and a
\textbf{Usability Report} in the winter.  These selections remain appropriate
for the project.  The interview report should document the target learner and
language-creator populations, interview methods, the questions asked, and
how stakeholder feedback changes the requirements.  The usability report should
use a systematic, quantified test plan with specific learner tasks, post-task
questions, and a survey for objective comparison against usability targets.

The team should not replace these CDs at this stage.  As implementation
progresses, the fall report and winter report plans will be linked together:
interview findings will inform the learner tasks and survey measures used in
the usability report.

\section{Team Charter Trigger Items}

The team charter establishes the following quantified triggers:
\begin{itemize}
    \item term attendance below 90\%, or two unexcused absences;
    \item an assigned task more than 48 hours past its due date without prior
    notice;
    \item no commits, code reviews, or closed issues for two consecutive weeks;
    and
    \item missing the internal 48-hour deadline for a deliverable section.
\end{itemize}

Based on the team and repository records available for the September 28 to
November 20 reporting period, no charter trigger violations are documented.
Attendance and task-timing records should continue to be captured in GitHub
meeting issues because the team has no complete set of attendance evidence in
the issue tracker yet.

If a trigger occurs, the charter's escalation process will be followed:
\begin{enumerate}
    \item hold a private check-in with the project lead;
    \item discuss the issue at a team meeting, record the decision and due dates
    in the meeting issue, and assign a specific recovery plan;
    \item reassign the member's tasks if the plan is not effective and reflect
    the situation honestly in the Team Contribution Survey; and
    \item escalate unresolved issues first to the TA and then to the
    instructor.
\end{enumerate}

The team will add meeting issues and attendance records, record internal
deliverable deadlines, and review the trigger data at every major milestone.
The 48-hour and two-week triggers will be retained unless the team finds them
impractical after using them with complete data.

\section{Additional Productivity Metrics}

\wss{If your team has additional metrics of productivity, please feel free to
add them to this report.  If not, please explicitly state that there are no additional
metrics.}

\wss{Additional metrics can include things like code reviews done, pull requests
created, count of joining meetings late, count of number of times contributions
had to be corrected, number of internal deadlines missed, test cases written, etc.}

\wss{We are looking for data on these metrics, not just a list of additional
metrics the team is planning on tracking.  However, if all you have is a plan,
please share it here.}

\end{document}